Parkinson's disease (PD) affects over 10 million individuals worldwide\cite{dorsey2018global}, with clinical presentation and progression trajectories varying substantially across patients\cite{fereshtehnejad2017subtypes,post2007progression}. While cross-sectional subtyping based on dominant motor phenotypes (tremor-dominant, postural instability-gait difficulty) has been established\cite{jankovic1990variable}, these classifications fail to capture longitudinal heterogeneity in progression rates—a critical determinant of disability accrual and treatment response\cite{latourelle2017large}. The inability to predict progression trajectories from baseline assessments limits personalized prognostication and impedes clinical trial design, where heterogeneous progression rates dilute treatment effect signals\cite{simuni2016predictors}.

Recent longitudinal cohort studies have revealed substantial variability in motor and cognitive decline: approximately 20-25\% of early-stage PD patients exhibit rapid motor deterioration (>5 UPDRS-III points/year), 50-60\% show moderate decline (2-4 points/year), and 15-25\% remain relatively stable (<2 points/year) over 5 years\cite{post2007progression,lawton2018developing}. However, these progression categories have been defined using arbitrary thresholds rather than data-driven clustering, and few studies have characterized their baseline predictors, biomarker profiles, or implications for trial enrichment. Existing subtyping efforts have largely focused on cross-sectional motor phenotypes\cite{fereshtehnejad2017subtypes} or genetic risk stratification\cite{nalls2019identification}, neglecting the longitudinal dimension essential for progression prediction.

Unsupervised machine learning offers a complementary approach: clustering longitudinal trajectories to discover latent subgroups without imposed clinical criteria\cite{schulam2019can,marquand2016conceptualizing}. However, trajectory clustering in PD faces methodological challenges. First, patients exhibit variable disease durations at baseline, necessitating \textit{latent time alignment} to a common disease progression timescale\cite{donohue2014preclinical}. Second, trajectories are high-dimensional and nonlinear, requiring dimensionality reduction while preserving temporal dynamics\cite{schulam2019can}. Third, progression subtypes must be validated beyond cluster validity metrics—demonstrating distinct baseline predictors, biological correlates, and clinical utility.

We address these challenges through a comprehensive trajectory clustering framework applied to the Parkinson's Progression Markers Initiative (\ppmi), a multicenter longitudinal study with standardized clinical, biomarker, and imaging assessments\cite{marek2011ppmi}. Our approach integrates variational autoencoder (VAE)-based trajectory modeling\cite{fortuin2020gp}, latent time alignment for disease duration normalization, and multi-resolution clustering with silhouette-based validation. We then characterize identified subtypes through baseline prediction modeling, biomarker profiling, and clinical trial simulations.

Our contributions are threefold. \textbf{First}, we identify three robust progression subtypes through data-driven clustering, validated across multiple metrics (silhouette scores, clinical distinctiveness, biomarker differences). \textbf{Second}, we demonstrate that subtype membership is partially predictable from baseline assessments (68\% accuracy), enabling prospective stratification. \textbf{Third}, we quantify clinical trial implications: enrichment with rapid progressors reduces sample size requirements by 38\% while improving effect size detectability, providing actionable guidance for adaptive trial designs. This framework establishes a roadmap for precision prognostication in PD and other neurodegenerative diseases.
